\section{Assumptions on Grammar and Semantic Correlation}
\label{sec:foundational_assumptions}

The averaging-based approximation developed in this chapter relies on a set of
assumptions about how grammatical and semantic information is distributed across
tokens in training data.
These assumptions are not specific to a particular model architecture, but rather
reflect generic statistical properties of natural language corpora and modern
training pipelines.
They are critical for ensuring that the proposed method isolates shared linguistic
structure while preserving the semantic correlations that must be learned.

\paragraph{Token indexing and notation.}
We index tokens by a pair $(j,t)$, where $j$ denotes the data sequence (e.g., a
sentence or document) and $t$ denotes the token position within that sequence.
Gradients, activations, and error signals are aggregated over tokens in a training
batch, which typically contains multiple sequences.

\paragraph{Assumption 1: Shared grammatical structure.}
We assume that all tokens, regardless of sequence or position, share a common set
of grammatical and structural regularities.
These include syntactic patterns, function words, positional conventions, and other
linguistic features that recur ubiquitously across natural language.
As a result, grammatical components of activations and gradients are \emph{globally
aligned} across tokens and sequences, and induce coherent, low-dimensional structure
in aggregated gradient matrices.

\paragraph{Assumption 2: Intra-sequence semantic correlation.}
Tokens within the same data sequence exhibit strong semantic correlation.
They share entities, topics, discourse context, and compositional meaning.
Consequently, semantic components of activations and gradients are correlated across
tokens belonging to the same sequence.
These correlations are essential for learning meaningful representations and must
not be suppressed by any approximation used during training.

\paragraph{Assumption 3: Inter-sequence semantic independence under data shuffling.}
We assume that the training data is well shuffled, so that different sequences in a
batch are drawn independently.
Under this assumption, semantic components of tokens belonging to \emph{different}
sequences have very low correlation in expectation.
While individual tokens may have large semantic gradients, their directions vary
across sequences and therefore cancel when aggregated across a sufficiently diverse
batch.

\paragraph{Extension to long-context training.}
In long-context settings, a single sequence may contain thousands of tokens.
Although tokens that are close in position are strongly correlated semantically,
semantic correlation decays with distance.
We therefore assume that tokens within the same sequence but separated by a
sufficiently large context window can be treated as approximately semantically
independent.
In this regime, averaging across such distant tokens behaves similarly to averaging
across independent sequences, while still preserving local semantic structure.

\paragraph{Implications for averaging-based approximation.}
Taken together, these assumptions imply a clear separation of roles:
\begin{itemize}
    \item Grammar-induced components are shared across tokens and sequences and
    reinforce coherently under aggregation, producing spiky, low-dimensional
    structure in gradient matrices.
    \item Semantic components are correlated within a sequence but largely
    uncorrelated across sequences (or across sufficiently distant positions in
    long-context settings).
    \item Averaging across tokens drawn from different sequences (or distant
    positions) suppresses independent semantic variation while preserving shared
    grammatical structure.
\end{itemize}

This separation is essential for the proposed approximation method.
It allows simple averaging operations to approximate the dominant, grammar-driven
low-rank structure of the agreement matrix, while leaving intra-sequence semantic
correlations intact.
As a result, the method isolates the spiky structure that dominates optimization
dynamics without interfering with the learning of semantic features.

\paragraph{Scope.}
Throughout this chapter, all averaging operations are implicitly understood to be
performed over tokens that are approximately semantically independent under the
assumptions above.
Violations of these assumptions---for example, highly correlated batches or
short-context averaging---may reduce the effectiveness of the approximation and are
outside the scope of the present analysis.